% The Quireloop guide. Every account gets its own copy as a project, made by
% server/src/lib/guide.js. Lines between %%%% PERSONAL and %%%% END are kept
% only on a personal install, lines between %%%% SERVER and %%%% END only on
% a lab server; the markers themselves are removed.
\documentclass[11pt]{article}
\usepackage[margin=2.2cm]{geometry}
\usepackage{graphicx}
\usepackage{xcolor}
\usepackage{booktabs}
\usepackage{array}
\usepackage[colorlinks=true,linkcolor=blue!50!black,urlcolor=blue!50!black]{hyperref}

\setlength{\parindent}{0pt}
\setlength{\parskip}{6pt}

% A screenshot, full width, with its caption.
\newcommand{\shot}[3][\linewidth]{%
  \begin{figure}[htbp]
    \centering
    \fbox{\includegraphics[width=\dimexpr#1-2\fboxsep-2\fboxrule\relax]{figures/#2}}
    \caption{#3}
  \end{figure}}
% A numbered marker, matching the red circles in the screenshots.
\newcommand{\mk}[1]{\textcolor{white}{\colorbox{red!75!black}{\textbf{#1}}}}
\newcommand{\key}[1]{\fbox{\footnotesize\texttt{#1}}}
\newcommand{\ui}[1]{\textbf{#1}}

\title{\textbf{Quireloop guide}\\[4pt]\large Everything Quireloop can do, step by step}
%%%% PERSONAL
\author{For Quireloop running on your own computer}
%%%% END
%%%% SERVER
\author{For Quireloop on your lab's server}
%%%% END
\date{}

\begin{document}
\maketitle

\begin{center}
\fbox{\begin{minipage}{0.9\linewidth}
\textbf{This project is yours to keep or delete.} It is an ordinary Quireloop
project, made for you when you first signed in. Read it here as a PDF, or open
\texttt{main.tex} to see how each page was written. Change anything you like:
it is a safe place to try every button. Deleting it does not bring it back.
\end{minipage}}
\end{center}

\tableofcontents
\newpage

%% ---------------------------------------------------------------------------
\section{What Quireloop is}

Quireloop is a LaTeX editor that runs in your web browser. You write LaTeX on
the left and see the PDF on the right. It saves every keystroke by itself, keeps
versions of your work, and lets several people edit the same paper at the same
time, each seeing the others' cursors.

%%%% PERSONAL
\textbf{You are using Quireloop on your own computer.} Your papers are files on
this computer, and nothing is sent anywhere unless you link it to a server
(Section~\ref{sec:server-link}). Quireloop compiles with the LaTeX installed on
this computer.
%%%% END
%%%% SERVER
\textbf{You are using Quireloop on your lab's server.} Your papers are stored on
the server, which keeps backups, and you can reach them from any browser.
Quireloop compiles with the LaTeX installed on the server, so you do not need
LaTeX on your own computer.
%%%% END

%% ---------------------------------------------------------------------------
\section{The dashboard: your papers}

After signing in you see the dashboard: every paper you own, and every paper
someone has shared with you.

%%%% PERSONAL
\shot{dashboard-personal.jpg}{The dashboard on your own computer.}

\begin{enumerate}
  \item[\mk{1}] Type a name for a new paper.
  \item[\mk{2}] Choose what it starts from: blank, an article, a report, slides,
    or a CV.
  \item[\mk{3}] \ui{New Project} creates it and opens it.
  \item[\mk{4}] \ui{Open a Folder} opens a folder that is already on this
    computer as a project. Quireloop edits the files where they are, so other
    programs keep seeing the same files.
  \item[\mk{5}] \ui{Import from Git} copies a paper from a git address.
  \item[\mk{6}] \ui{Upload Project (.zip)} opens a zipped LaTeX project. Quireloop
    finds the main file by itself.
  \item[\mk{7}] \ui{Download all my papers} saves one \texttt{.zip} of everything.
  \item[\mk{8}] \ui{Work with a server} links this computer to a lab server, so
    your server papers are copied here and kept in step
    (Section~\ref{sec:server-link}).
  \item[\mk{9}] \ui{Keep a backup copy in a folder} copies your papers as
    ordinary files to a folder you choose, every few minutes
    (Section~\ref{sec:backup}).
  \item[\mk{10}] Your account: password, two-step sign-in and the assistant's
    settings (Section~\ref{sec:account}).
\end{enumerate}
%%%% END
%%%% SERVER
\shot{dashboard-server.jpg}{The dashboard on a lab server.}

\begin{enumerate}
  \item[\mk{1}] Type a name for a new paper, choose a starting template next to
    it, and press \ui{New Project}. You can also open a folder you have on the
    server, import from git, or upload a \texttt{.zip}.
  \item[\mk{2}] Each paper is a card. Click its name to open it. Papers shared
    with you are marked \ui{Shared}.
  \item[\mk{3}] \ui{Download all my papers} saves one \texttt{.zip} of every paper
    you own or share. Do this now and then, and keep the file somewhere safe.
  \item[\mk{4}] \ui{Admin} appears only for administrators
    (Section~\ref{sec:admin}).
  \item[\mk{5}] Your account: password, two-step sign-in and the assistant's
    settings (Section~\ref{sec:account}).
\end{enumerate}
%%%% END

%% ---------------------------------------------------------------------------
\section{The project screen}

Opening a paper shows three parts side by side: your files, the editor and the
PDF. Quireloop compiles the paper as soon as you open it, so the PDF is already
there.

\shot{project.jpg}{A paper open in Quireloop.}

\begin{enumerate}
  \item[\mk{1}] The side panel. \ui{Files} lists the project's files and folders;
    \ui{Outline}, \ui{Git} and \ui{Search} are the other tabs.
  \item[\mk{2}] The \textbf{insert toolbar}: formatting, structure, maths,
    tables, figures and packages in one click (Section~\ref{sec:toolbar}).
  \item[\mk{3}] The editor. Type LaTeX here. It saves by itself; there is no save
    button.
  \item[\mk{4}] \ui{Compile} builds the PDF. Shortcut: \key{Ctrl+S} or
    \key{Ctrl+Enter} (\key{Cmd} on a Mac).
  \item[\mk{5}] The PDF. Zoom, find text, and double-click to jump to the source
    (Section~\ref{sec:sync}).
  \item[\mk{6}] \ui{Share} lets other people in (Section~\ref{sec:share}).
  \item[\mk{7}] \ui{History} lists versions of the paper
    (Section~\ref{sec:history}).
  \item[\mk{8}] \ui{Menu}: compiler, auto-compile, clean-up and download.
  \item[\mk{9}] \ui{Edit}: the main file, writing tools and collaboration.
\end{enumerate}

Panels that open from the top bar (Share, History, menus) close when you click
anywhere else, or press \key{Esc}. Drag the borders between the three parts to
resize them; Quireloop remembers the sizes.

%% ---------------------------------------------------------------------------
\section{Writing faster with the insert toolbar}\label{sec:toolbar}

The row above the editor holds the commands used most often. Select some text
first to wrap it, or press a button with nothing selected to insert a
placeholder you can type over.

\shot{insert-toolbar.jpg}{The insert toolbar, in five groups.}

\begin{enumerate}
  \item[\mk{1}] \textbf{Text:} \ui{B} bold, \ui{I} italic, \ui{U} underline,
    \ui{H} highlight, \ui{Colour} coloured text, and \ui{Align} for left,
    centred or right-aligned text.
  \item[\mk{2}] \textbf{Structure:} \ui{Heading} (section, subsection,
    subsubsection, paragraph), bulleted and numbered lists, \ui{fn} footnote,
    \ui{[1]} citation, \ui{ref} reference and \ui{lbl} label.
  \item[\mk{3}] \textbf{Maths:} \ui{\$x\$} inline maths, \ui{(1)} a numbered
    equation, \ui{a=} aligned equations, \ui{½} a fraction, \ui{[:]} a matrix,
    and \ui{Symbol} for Greek letters, operators and arrows.
  \item[\mk{4}] \textbf{Tables, figures and packages:} \ui{Table} asks for rows
    and columns and builds the whole table; \ui{Figure} lists the images in your
    project; \ui{+ Package} adds a package to the preamble, with a short note on
    what each common package is for.
  \item[\mk{5}] \textbf{Text size:} \ui{A$-$} and \ui{A+} make the editor's text
    smaller or larger; the percentage resets it. Each browser remembers its
    own size.
\end{enumerate}

\textbf{Packages are added for you.} Highlighting needs the \texttt{soul} and
\texttt{xcolor} packages, a table needs \texttt{booktabs}, and so on. When a
button needs one, Quireloop adds the \verb|\usepackage| line to the preamble
and shows a short note saying so. In a chapter file without a preamble, the
line goes into your main file.

\subsection*{More writing help}
\begin{itemize}
  \item \textbf{Autocomplete:} start typing a command, such as \verb|\beg|, and
    choose from the list. \verb|\cite{| offers every key in your
    \texttt{.bib} files.
  \item \textbf{Hover previews:} point at maths to see it typeset, or at a
    citation key to see the reference.
  \item \textbf{Paste or drop images:} paste a screenshot or drag an image into
    the editor. It is saved in \texttt{figures/} and a figure is inserted.
  \item \textbf{Folding:} the small arrows next to the line numbers fold a
    section or environment away.
\end{itemize}

%% ---------------------------------------------------------------------------
\section{Menu and Edit}

\shot[0.75\linewidth]{menu.jpg}{The Menu.}

\begin{enumerate}
  \item[\mk{1}] \ui{Auto-compile} recompiles shortly after you stop typing.
  \item[\mk{2}] \ui{Compiler}: pdfLaTeX for most papers; XeLaTeX or LuaLaTeX for
    system fonts and \texttt{fontspec}.
  \item[\mk{3}] \ui{Clean aux files} removes LaTeX's leftover files. Use it when
    a compile stays broken after you have fixed the mistake.
  \item[\mk{4}] \ui{Download .zip} saves this paper's files.
\end{enumerate}

\shot[0.75\linewidth]{edit-menu.jpg}{The Edit menu.}

\begin{enumerate}
  \item[\mk{1}] \ui{Main file}: the \texttt{.tex} file that is compiled into the
    PDF. A paper split into chapters has one main file that pulls the others in
    with \verb|\input| or \verb|\include|. Choose it here; the PDF is rebuilt
    straight away.
  \item[\mk{2}] \ui{Spell check}, and \ui{Grammar check} when your Quireloop
    has it set up. \ui{Copy text without LaTeX} copies plain prose to paste into
    other writing tools.
  \item[\mk{3}] \ui{Insert symbol} opens a palette of symbols and snippets.
  \item[\mk{4}] \ui{Suggest mode} turns your edits into suggestions that others
    accept or reject (Section~\ref{sec:together}).
  \item[\mk{5}] \ui{Comments} on selected text.
  \item[\mk{6}] \ui{Chat} with everyone working on the paper.
  \item[\mk{7}] \ui{Assistant}: ask an AI model about your paper
    (Section~\ref{sec:assistant}).
\end{enumerate}

%% ---------------------------------------------------------------------------
\section{Compiling and fixing mistakes}

Press \ui{Compile} (or \key{Ctrl+S}). Quireloop runs LaTeX, BibTeX or Biber as
many times as needed, and shows the PDF even when LaTeX reports errors, as long
as it produced one.

\shot{logs.jpg}{A compile with a mistake.}

\begin{enumerate}
  \item[\mk{1}] \ui{Logs} shows a small badge with the number of problems.
    Press it to open the list.
  \item[\mk{2}] Each problem names its file and line. Click it to jump there; the
    line is also marked in the editor with a red dot and a wavy underline.
  \item[\mk{3}] \ui{Show raw log} shows LaTeX's own output, for the rare
    message the list does not explain.
\end{enumerate}

\subsection{Jumping between the PDF and the source}\label{sec:sync}
\begin{itemize}
  \item \textbf{Source to PDF:} double-click a word in the editor, or press
    \key{Ctrl+Alt+J}. The PDF scrolls to it and outlines it.
  \item \textbf{PDF to source:} double-click a word in the PDF. The editor opens
    the right file and selects the word.
\end{itemize}
In the PDF pane you can also select and copy text, search with
\ui{Find in PDF}, and switch to a dark page view with the half-moon button.

%% ---------------------------------------------------------------------------
\section{Files, folders and the outline}

In the \ui{Files} tab, \ui{New file} creates a file (type
\texttt{chapters/intro.tex} to create it inside a folder), \ui{Folder} creates a
folder, and \ui{Upload} adds files from your computer. Each file has buttons to
rename and delete it. Folders start closed; click one to open it.

\shot[0.4\linewidth]{outline.jpg}{The outline of a paper.}

The \ui{Outline} tab lists the open file's sections. They start folded, so a
long paper shows only its main sections.
\begin{enumerate}
  \item[\mk{1}] The arrow opens or folds a section's subsections.
  \item[\mk{2}] Click any title to jump to it in the editor.
  \item[\mk{3}] \ui{Open all} or \ui{Fold all} at once.
\end{enumerate}

The \ui{Search} tab searches every file in the project. To find a command or a
file by name, press \key{Ctrl+K} (\key{Cmd+K} on a Mac):

\shot[0.7\linewidth]{palette.jpg}{Search for any command or file.}

%% ---------------------------------------------------------------------------
\section{Sharing a paper}\label{sec:share}

Press \ui{Share} next to \ui{Compile}. Only the paper's owner can change who has
access.

\shot[0.55\linewidth]{share.jpg}{The Share panel.}

\begin{enumerate}
  \item[\mk{1}] \ui{Add people} lists everyone else with an account here, so you
    do not need to know their username or email. Tick the people to add.
  \item[\mk{2}] Choose a role for each: an \textbf{Editor} can change the paper;
    a \textbf{Viewer} can read, compile and comment, but not change anything.
  \item[\mk{3}] Press \ui{Share with \dots} to give them access. The paper
    appears on their dashboard at once.
  \item[\mk{4}] The people who already have access. Change a role or press
    \ui{Remove} at any time.
  \item[\mk{5}] A \textbf{share link} lets in anyone with an account who opens
    it, with the role you choose. It suits someone who has no account yet: they
    sign up first, then open the link.
  \item[\mk{6}] \ui{Create link}, then \ui{Copy} it. \ui{Revoke} switches a link
    off.
\end{enumerate}

%% ---------------------------------------------------------------------------
\section{Working together}\label{sec:together}

Everyone with access can edit at the same time. Each person's cursor appears in
their own colour, with their name. There is no locking and no ``someone else
saved first'': edits are merged as you type. The dot next to the paper's name is
green when you are connected; if your connection drops, keep typing, and your
work is merged when it returns.

\begin{itemize}
  \item \textbf{Comments:} select text, then \ui{Edit} $\rightarrow$
    \ui{Comment on selection}. Reply, resolve or delete in the Comments panel.
    A comment stays with its text even as the paper changes around it.
  \item \textbf{Suggest mode:} with \ui{Suggest mode} on, what you type is shown
    in green and what you delete is struck through, without being removed.
    Anyone who can edit accepts or rejects each suggestion in the Suggestions
    panel.
  \item \textbf{Chat:} a message board for the paper, under \ui{Edit}.
\end{itemize}

%% ---------------------------------------------------------------------------
\section{Version history}\label{sec:history}

Press \ui{History}. Quireloop keeps three kinds of versions:
\begin{itemize}
  \item \textbf{Autosave (writing only):} every minute in which the writing
    changed, whether you compiled or not, and whether the change was made in
    Quireloop or by another program. Holds your text files only.
  \item \textbf{Auto (compile):} a full copy after each successful compile.
  \item \textbf{Saved:} one you save yourself with \ui{Save current version}, at
    the top of the panel.
\end{itemize}

\shot[0.5\linewidth]{history.jpg}{Version history.}

\begin{enumerate}
  \item[\mk{2}] A version you named. Name any version, an autosave included, to
    find it again later: ``sent to co-authors'', ``before the rewrite''.
  \item[\mk{3}] \ui{Diff} shows what changed between that version and now.
  \item[\mk{4}] \ui{Name} or \ui{Rename}.
  \item[\mk{5}] \ui{Restore} brings the paper back to that version. The current
    state is saved first, as \ui{Before restore}, so a restore can always be
    undone. Everyone with the paper open sees the restored text straight away.
    Restoring an autosave puts back the writing and leaves figures as they are.
\end{enumerate}

%% ---------------------------------------------------------------------------
\section{Editing with other programs}

You can change a paper's files with any other program while Quireloop is
running: another editor, a script, or an AI assistant that writes files. Even
with the file open in Quireloop, the change appears in the editor within a few
seconds and is never written over. If someone types in the browser at the same
moment, both changes are kept.

%%%% PERSONAL
On this computer, a project's files are in Quireloop's data folder, or wherever
the folder you opened with \ui{Open a Folder} is.
%%%% END
%%%% SERVER
On the server, this applies to programs running on the server itself, such as a
folder you opened with \ui{Open a Folder} and edit over a remote session.
%%%% END

%% ---------------------------------------------------------------------------
\section{Git}

Every project is also a git repository, in the \ui{Git} tab.

\shot[0.32\linewidth]{git.jpg}{The Git tab.}

Write a message and press \ui{Commit} to record the changes listed below it.
\ui{Set remote} connects the project to a repository on a git host, using its
address and an access token; then \ui{Pull} and \ui{Push} exchange changes.
Quireloop's own files (the build folder, versions, comments) are never pushed.

%% ---------------------------------------------------------------------------
\section{The writing assistant}\label{sec:assistant}

\ui{Edit} $\rightarrow$ \ui{Assistant} opens a panel where you can ask about the
open file, or about selected text: fix a LaTeX error, shorten a paragraph, draft
a caption. Nothing changes in your paper until you press \ui{Insert}, and
\key{Ctrl+Z} undoes it.

Choose where answers come from under your account's \ui{AI Assistant} settings:
local models on a server (free, and your text stays there), or your own API key
from an AI provider. Only you can use your key.

%% ---------------------------------------------------------------------------
\section{Your account}\label{sec:account}

Click your email at the top of the dashboard.
\begin{itemize}
%%%% SERVER
  \item If you sign in with your server account, your password and two-step
    sign-in are those of the server, and are changed there.
%%%% END
  \item \textbf{Change password:} needs your current password.
  \item \textbf{Two-step sign-in (2FA):} scan the code with an authenticator
    app, then enter the six-digit number it shows. Every sign-in then asks for a
    current number. Keep the app safe: losing it locks you out until an admin
    helps.
  \item \textbf{AI Assistant:} where the assistant's answers come from.
\end{itemize}

%%%% PERSONAL
%% ---------------------------------------------------------------------------
\section{On your own computer}

\subsection{Starting Quireloop by itself}
Quireloop can start when you log in to your computer and restart if it ever
stops. In a terminal, in the folder where Quireloop is installed, run
\begin{center}\texttt{npm run service:install}\end{center}
This works on macOS, Windows and Linux, and needs no administrator rights.
To undo it, run \texttt{npm run service:uninstall}. If you installed Quireloop
with \texttt{pip}, start it with \texttt{quireloop start} instead.

\subsection{Linking to a lab server}\label{sec:server-link}
Under \ui{Work with a server} on the dashboard, enter the server's address and a
name for this computer, and press \ui{Connect}. Sign in to the server on the page
that opens and approve the link. From then on:
\begin{itemize}
  \item every paper you can open on the server is copied to this computer, and
    kept in step every 30 seconds while you are online;
  \item you always work on the copy here, so losing the internet changes
    nothing; your edits and the server's are merged when it returns;
  \item a file you have open is linked live, so you see co-authors' typing as it
    happens.
\end{itemize}
Comments, chat, suggestions and version history stay on each side and are not
copied.

\subsection{A backup folder}\label{sec:backup}
Under \ui{Keep a backup copy in a folder}, enter a folder and press \ui{Save}.
Your papers are copied there as ordinary files at once, then every few minutes.
Nothing in that folder is ever deleted: when a file changes, its earlier copy
moves to \texttt{\_older-versions}. Choose a folder that your cloud storage app
syncs to keep a copy off this computer too.

\subsection{Compiling here}
Compiling uses the LaTeX installed on this computer. A paper that needs a
package you do not have will report it in \ui{Logs}; install the package with
your LaTeX distribution's package manager.
%%%% END

%%%% SERVER
%% ---------------------------------------------------------------------------
\section{On the lab server}

\begin{itemize}
  \item \textbf{Where your papers are:} on the server, backed up regularly to a
    separate disk. Still, press \ui{Download all my papers} now and then.
  \item \textbf{Working offline:} install Quireloop on your own computer and link
    it to this server; your papers are then copied to your computer and kept in
    step, and you can work without internet.
  \item \textbf{Opening a folder:} \ui{Open a Folder} on the dashboard lists the
    folders on the server you are allowed to open, such as your home folder.
    Quireloop edits those files where they are.
  \item \textbf{Compile limits:} each compile runs separately from everyone
    else's and stops after a few minutes, so a runaway paper cannot slow the
    server for others.
\end{itemize}

\subsection{For administrators}\label{sec:admin}
The \ui{Admin} button on the dashboard is shown only to administrators.
\begin{itemize}
  \item \textbf{Invites:} create a single-use invite code or link for each new
    person.
  \item \textbf{Users:} disable an account to lock it out at once (its papers are
    kept), or enable it again; assign a member's folder on the server.
  \item \textbf{Backups and settings:} the time of the last backup, and the
    server-wide assistant settings.
\end{itemize}
Administrators cannot open other people's papers from the admin panel; access
always goes through sharing.
%%%% END

%% ---------------------------------------------------------------------------
\section{Keyboard shortcuts}

On a Mac, use \key{Cmd} where this table says \key{Ctrl}.

\begin{center}
\begin{tabular}{>{\raggedright}p{0.35\linewidth}p{0.5\linewidth}}
\toprule
\textbf{Keys} & \textbf{What they do} \\
\midrule
\key{Ctrl+S}, \key{Ctrl+Enter} & Compile \\
\key{Ctrl+K} & Search commands and files \\
\key{Ctrl+Alt+J} & Show the cursor's place in the PDF \\
\key{Ctrl+Z}, \key{Ctrl+Y} & Undo and redo your own edits \\
\key{Ctrl+F} & Find in the editor, or in the PDF when it has focus \\
\key{Esc} & Close a panel or menu \\
Double-click & Jump between the source and the PDF \\
\bottomrule
\end{tabular}
\end{center}

%% ---------------------------------------------------------------------------
\section{When something goes wrong}

\begin{description}
  \item[The PDF does not change.] Open \ui{Logs}: an error may have stopped
    the compile, and the last good PDF is still shown. If the log looks fine,
    use \ui{Menu} $\rightarrow$ \ui{Clean aux files} and compile again.
  \item[The wrong file is compiled.] Choose the right one under \ui{Edit}
    $\rightarrow$ \ui{Main file}.
  \item[I need an older version.] \ui{History}, then \ui{Diff} to check and
    \ui{Restore} to bring it back.
  \item[The dot next to the paper's name is not green.] You are offline. Keep
    working; your edits are kept and merged when the connection returns.
  \item[A package is missing.] The \ui{Logs} list names it.
%%%% PERSONAL
    Install it with your LaTeX distribution's package manager.
%%%% END
%%%% SERVER
    Ask your administrator to install it on the server.
%%%% END
  \item[``Too many attempts'' when signing in.] Wait for the time shown, then
    try again.
\end{description}

\end{document}
